\begin{tabular}{lllll}
\toprule
\textbf{Scenario} & \textbf{Contrast} & \textbf{\ensuremath{\Delta} average fit [95\% CI]} & \textbf{\ensuremath{\Delta} fit floor [95\% CI]} & \textbf{\ensuremath{\Delta} orientation SD [95\% CI]} \\
\midrule
Normal load & COCPIT only \ensuremath{-} Baseline & +0.2 [0.0, 0.5] & +0.4 [0.0, 0.8] & -0.09 [-0.27, 0.09] \\
Normal load & NURSTA only \ensuremath{-} Baseline & -0.1 [-0.4, 0.2] & -0.1 [-0.5, 0.3] & +0.07 [-0.10, 0.23] \\
Normal load & Both \ensuremath{-} Baseline & \textbf{+0.3 [0.1, 0.5]} & +0.4 [0.0, 0.8] & -0.13 [-0.34, 0.08] \\
\addlinespace[2pt]
High load & COCPIT only \ensuremath{-} Baseline & \textbf{+0.4 [0.1, 0.7]} & +0.7 [0.0, 1.4] & \textbf{-0.25 [-0.48, -0.01]} \\
High load & NURSTA only \ensuremath{-} Baseline & +0.0 [-0.3, 0.3] & +0.2 [-0.4, 0.8] & \textbf{-0.22 [-0.43, -0.02]} \\
High load & Both \ensuremath{-} Baseline & \textbf{+0.5 [0.2, 0.7]} & \textbf{+0.8 [0.3, 1.3]} & \textbf{-0.27 [-0.44, -0.11]} \\
\bottomrule
\end{tabular}

\par\footnotesize\textit{Note.} The fit floor is the lowest of the five orientation scores within each seed. Orientation SD is the within-seed standard deviation across the five designed orientations; a negative change indicates less disparity. Confidence intervals use seed-paired means (n=10). These are synthetic-orientation ensemble metrics, not population accessibility estimates. Bold values have 95\% confidence intervals that exclude zero.
